\subsection{Governance}

Article 36 of Protocol I to the Geneva Conventions requires states to conduct legal reviews of new weapons systems to determine whether their use would be prohibited under international law~\citep{icrc_api_article36_1977}. Israel is among the few states that have not ratified Protocol I~\citep{icrc_api_stateparties}, meaning it has no legal obligation to assess Habsora's compliance with IHL prior to deployment. UNESCO further holds that delegating life-or-death decisions to an AI system is itself impermissible~\citep{aimag.v40i1.2848}, yet without ratification, no mechanism exists to enforce this in the Israeli context.

Once deployed, accountability becomes severely diffuse. Under IHL, commanders remain legally responsible for how weapons are used, even when autonomous functions are involved~\citep{icrc_api_article86_1977}, and are judged on what they knew at the time of the decision~\citep{henckaerts2005customary}. In practice, responsibility is spread across developers, the algorithm, the receiving IDF officer, and the approving commander~\citep{aimag.v40i1.2848}. The system's documented 10\% error tolerance in targeting recommendations illustrates the stakes: at operational scale, this margin translates into foreseeable civilian casualties, yet no single actor bears clear responsibility. In high-tempo environments, real-time monitoring is not feasible, compounding this further.

Post-deployment review processes are largely absent, and the system's opacity makes it difficult to determine whether casualties resulted from an AI recommendation or a human override. Civilians — the most directly affected stakeholders — were never consulted, and no structured ethical assessment, such as Value Sensitive Design, was conducted. This absence of accountability and stakeholder engagement carries implications not only for this conflict, but for the global development of IHL norms.

\subsection{Socio-technical}

Palestinian civilians, the stakeholders most directly affected by Habsora, were not involved in its design and their values were not elicited to ground the system's targeting logic. They remain non-consenting parties who bear the consequences of the system's use without recourse, reflecting a profound asymmetry between those operating the system and those subjected to it.

For IDF commanders, the conditions necessary for genuine human oversight are structurally absent. \citet{mhcof} argue that meaningful human control requires both adequate information and meaningful freedom of action. Habsora undermines both: commanders receive recommendations without transparency into the underlying reasoning, and the short decision windows afforded by the system leave little room for critical evaluation. \citet{miller2018explanationartificialintelligenceinsights} similarly demonstrates that without understanding what went into a decision, the explainee cannot meaningfully interrogate it. Combined with the sheer volume of targets the system generates, this reduces what would ordinarily involve extended deliberation and proportionality assessments to a button press. Former IDF commanders have described this process as rubber-stamping~\citep{kozlovski2024algorithms}, corroborating the theoretical prediction of automation bias, whereby operators over-rely on AI outputs under time pressure~\citep{parasuraman2010complacency}.

The consequence is that IDF commanders formally bear moral and legal responsibility for strikes they cannot meaningfully evaluate, while the system's opacity diffuses accountability in practice. Targeting decisions that once demanded careful human judgement are now effectively delegated to an algorithm, hollowing out the very human control that IHL presupposes and that \citet{mhcof} identify as a precondition for attributing responsibility.